\section{误差分析、模型检验与灵敏度}

各问的结果已在相应章节给出局部敏感性。本节做三项横向汇总：数值正确性的核验、误差来源的分解、以及结论对关键假设的稳健性。

\subsection{数值正确性核验}

表~\ref{tab:verify} 汇总了本文为各问设置的正确性检验。所有检验均以独立于求解路径的方式实施（解析解、有限差分、资源守恒），而非重复同一次计算。

\begin{table}[htbp]
\centering
\caption{数值正确性核验汇总}
\label{tab:verify}
\small
\begin{tabular}{l|l|l|r}
\toprule
\textbf{问题} & \textbf{检验方式} & \textbf{检验内容} & \textbf{结果} \\
\midrule
一 & 一对一连接核对 & A4/A5 等六对表的 \texttt{index} 与列序 & 全部一致 \\
一 & 归一化核对 & 配比行和偏离 1 的最大值 & $<0.4\%$，已行归一化 \\
一 & bootstrap 稳定性 & 1768 个交互系数的 $|b|/\mathrm{sd}$ & 27.3\% 显著 \\
二 & 有限差分 & 解析弹性 vs 中心差分 & 偏差 $<1.8\times10^{-12}$ \\
二 & 解析参照 & 分组留出 vs 全样本拟合 & 参数一致 \\
三 & 解析参照 & 式\eqref{eq:q3_analytic} vs 数值搜索（15 组） & 相对偏差 $<1.1\times10^{-6}$ \\
三 & 资源守恒 & 三项成本之和 vs 预算 $C$（45 组） & 相对残差 $\le2.2\times10^{-16}$ \\
三 & 一阶条件 & $|\partial L/\partial\ln N|$ 与 $Q$ 的 KKT 符号 & 全部满足 \\
四 & 聚合口径 & BBH 宏平均 vs 题量加权 vs 榜单值 & 仿射关系 $\rho=0.9979$ \\
四 & 分解加和 & 规模 $+$ 非规模 $+$ 残差 $=\Delta S$ & 恒等成立 \\
四 & 无时间泄漏 & 回测仅用截止点之前的数据 & 已核验 \\
\bottomrule
\end{tabular}
\end{table}

其中问题三的解析参照检验值得单独说明：本文最初使用常规二维 Nelder--Mead 搜索时，数值解与解析解相差 60\%，但损失差仅 $9.5\times10^{-6}$，极易被误判为``已经收敛''。正是解析参照暴露了这一问题，促使本文改用``外层网格 $+$ 内层一维精确极小化''的策略，修正后偏差降至 $10^{-6}$ 量级。这说明\textbf{在平坦损失面上，只看损失值不足以判断解的正确性}，必须与解析参照对照。

\subsection{误差来源分解}

本文的误差可归为四类，表~\ref{tab:errsrc} 给出各自的量级与影响范围。

\begin{table}[htbp]
\centering
\caption{误差来源与影响}
\label{tab:errsrc}
\small
\begin{tabular}{l|l|r|l}
\toprule
\textbf{误差来源} & \textbf{受影响结论} & \textbf{实测量级} & \textbf{处理方式} \\
\midrule
数据来源虚标（B1 为生成数据） & 标度律参数的泛化性 & 相对残差 $5.9\times10^{-5}$ & 改用 B4/B5 作泛化检验 \\
跨源坐标不可比（B2） & 族外验证 & 平移 1.18 nat & 只做排序验证 \\
质量域映射代理（11 个域） & $Q(p)$ 的绝对值 & 未直接测得 & 只用于单调折算，不做因果解读 \\
桥接误差（Loss$\to$Score） & 问题四绝对水平 & 留一 MAE 7.34 分 & 计入预测区间；只作趋势解读 \\
损失面平坦 & 问题三的 $N^\ast$ & 辨识带 2.3$\sim$91 倍 & 报告区间而非点估计 \\
组成变化 & 问题四贡献分解 & 前沿口径残差 62\% & 主用固定组成口径 \\
\bottomrule
\end{tabular}
\end{table}

\subsection{关键假设的稳健性}

\textbf{假设 3（有效数据量形式）}：改用双通道模型 $A(NQ^\rho)^{-\alpha}$ 后，$\rho$ 与 $\kappa$ 不可辨识，置信区间互相覆盖。这说明现有数据无法区分``质量作用于数据效率''与``质量作用于参数效率''，但\textbf{两种形式在 $Q\in[0.5,1]$ 内的预测差异小于 $10^{-3}$ nat}，对问题三的配置结论无实质影响。

\textbf{假设 4（预算紧约束）}：已在 45 个情景上逐点核验，成本残差为机器精度，未发现约束松弛的情形。

\textbf{假设 6（跨源同尺度）}：B6 在 $Q=1$ 的 45 个点相对 B1 经典律的残差均值为 $+0.0036$、标准差 0.055。若把基准 $Q$ 从 1 改为 0.95（即假设 B6 的``完美质量''并非真正完美），$\kappa$ 的估计仅变化约 3\%，结论稳健。

\textbf{假设 8（开源口径）}：严格与宽松口径的前沿基准相差 1.44 分（36.26 对 37.70），远小于 12 个月的预测增量（约 12.9 分），故口径选择不改变主要结论。

\textbf{质量成本函数}：三类函数的启动点落在 $\log_{10}C\in[19.6,21.1]$，跨度 1.5 个数量级。这意味着``质量投入存在启动预算''是稳健结论，但\textbf{启动点的具体位置依赖于成本函数的选择}，不应给出单一数值。

\subsection{与基线模型的对照}

为说明本文模型相对简单做法的增益，表~\ref{tab:baseline} 给出与``最低参照''的对照。

\begin{table}[htbp]
\centering
\caption{与基线做法的对照}
\label{tab:baseline}
\small
\begin{tabular}{l|l|r|l}
\toprule
\textbf{问题} & \textbf{基线做法} & \textbf{本文结果} & \textbf{增益} \\
\midrule
一（配比） & 训练集均值预测 & CV RMSE 0.257 & 基线 RMSE 0.345，改善 25.5\% \\
一（配比） & 线性混料 & CV RMSE 0.257 & 线性为 0.283，改善 9.2\% \\
一（质量） & 等权算术平均 & 秩相关 0.665 & 分组几何平均携带额外信息 \\
二（质量） & 不含 $Q$ 的经典律 & RMSE 0.117 & 常数基准 0.345，改善 66.1\% \\
二（质量） & $\kappa=1$ 的约束模型 & 秩相关 $>0.999$ & 与自由估计几乎等价 \\
四（能力） & 仅用参数量 & $R^2$ 0.465 & 加入时间与类型效应 \\
四（分解） & 季度均值分解 & 固定组成分解 & 避免组成变化造成的 $-122\%$ 伪贡献 \\
\bottomrule
\end{tabular}
\end{table}

最后一行值得展开：若直接用季度\textbf{均值}做分解，会得到规模贡献 $-122\%$、非规模贡献 $+1974\%$ 这类超出合理范围的数值。这不是模型的错误，而是\textbf{组成变化的污染}——提交的模型越来越小，使``平均参数量''下降。改用固定组成子集后，未解释部分从 62\% 降到 8.5\%。这说明\textbf{在做贡献分解前必须先控制组成变化}，否则会得出``规模扩张拖累了能力''的错误结论。
